%%%%%%%%%%%%%%%%%%%%%%%%%%%%%%%%%%%%%%%%%%%%%%%%%%%%%%%%%%%%%%%%%%%%%%
% Lớp 12 | Ứng dụng đạo hàm để khảo sát hàm số | Khảo sát sự biến thiên và vẽ đồ thị hàm số
% Dạng 2D1?5-6: Phương trình tiếp tuyến của đồ thị hàm số
% Số câu: 7  (sinh tự động bởi tach_cau.py; từ nay thầy sửa trực tiếp file này)
%%%%%%%%%%%%%%%%%%%%%%%%%%%%%%%%%%%%%%%%%%%%%%%%%%%%%%%%%%%%%%%%%%%%%%

%%%%%%%%%%%%%%%% TRẮC NGHIỆM 4 PHƯƠNG ÁN %%%%%%%%%%%%%%%%

% nguồn: Phan1_Lop_12_Dai_so_Chuong_1
% nguồn: Phan1_Lop_12_Dai_so_Chuong_1

%%%%%%%%%%%%%%%% ĐÚNG/SAI %%%%%%%%%%%%%%%%

% nguồn: Phan1_Lop_12_Dai_so_Chuong_1
% nguồn: Phan1_Lop_12_Dai_so_Chuong_1

%%%%%%%%%%%%%%%% TRẢ LỜI NGẮN %%%%%%%%%%%%%%%%

% nguồn: Phan2_Lop_12_Dai_so_Chuong_1
% nguồn: Phan2_Lop_12_Dai_so_Chuong_1
\begin{bt}%[2D1V5-6]%[ma:c243bd5d]
	\immini{Một đường ray tàu lượn trong khu vui chơi giải trí có hình dáng được  mô phỏng theo đồ thị của hàm số $y=x^3-3x+2$, ký hiệu là $(C)$. Để đảm bảo an toàn và tính thẩm mỹ, người ta chọn hai điểm $A(a;b)$ và $B(c;d)$ trên đường ray sao cho tiếp tuyến tại hai điểm này có cùng độ dốc (cùng hệ số góc). Đồng thời, đoạn đường nối hai trụ đỡ tại các điểm $A$ và $B$ phải vuông góc với một
	đường dây điện có phương trình $x+y-5=0$.  Tìm $b+d$.}{
	\begin{tikzpicture}[line cap=butt,line join=miter,>=stealth]
		\draw[->] (-3,0)--(3,0) 
			 node[shift={(-100:7pt)},font=\normalsize]{$x$};
		\draw[->]  (0,-1)--(0,5)
			node[shift={(180:7pt)},font=\normalsize]{$y$}; 
		\draw (5pt,0) |- (0,5pt) 
			(0,0)  node[shift={(225:9pt)},font=\normalsize]{$O$};
		\fill (0,0)  circle (1pt);
		\tikzset{declare  function={f (\x)=(\x)^3-3*(\x)+2;}}
		\begin{scope}
			\draw[blue,samples=100] plot[domain=-2.1:1.8] (\x,{f  (\x)});
		\end{scope}
	\end{tikzpicture}
	 }
	 \shortans{4}
	 \loigiai{Ta có $y=x^{3}-3x+2\Rightarrow y'=3x^{2}-3$.\\
		Tiếp tuyến với $(C)$ tại $A$, $B$ có cùng hệ số góc và chỉ khi 
		 \[f'(x_A)=f'(x_B)
		\Leftrightarrow  x_A^2=x_B^2
		\Leftrightarrow \hoac{&x_A=x_B\,(\text{loại}) \\&x_A+x_B=0.}\] 
		Suy ra $A$, $B$ đối xứng với nhau qua điểm $I(0;2)$ là tâm đối xứng của  $(C)$.\\
		$AB\perp  d\colon x+y-5=0\Rightarrow AB\colon x-y+m=0$.\\
		$AB$ qua  $I$ nên ta có $m=2\Rightarrow AB\colon x-y+2=0$.\\
		Khi đó hoành độ  $A$, $B$ thỏa mãn phương trình
		$x^{3}-3x+2=x+2\Leftrightarrow \hoac{&x=0\,(\text{loại})  \\
		 &x=\pm 2.}$\\
		$\Rightarrow  A(2;4)$, $B(-2;0)$.\\
		Do đó, $b=4$  và $d=0$.\\
		Vậy $b+d=4+0=4$. 
	 }
\end{bt}


%%%%%%%%%%%%%%%% TỰ LUẬN %%%%%%%%%%%%%%%%

% nguồn: Phan2_Lop_12_Dai_so_Chuong_1
% nguồn: 2026-11-16_12_2D1_TL_Ham_so (nhập 16/11/2026)
\begin{bt}%[2D1H5-6]%[ma:84b59bb1]
% nguon: Archive/2024/TaiLieuGiangDay/1.HamSo/1. K12 - 4in1 - Hamso/ThubaiGoc/Bai5-Dang3.tex #2140
[Mức độ 2]%[Bài giảng 12 mới, Lê Hồng Phi]
Cho hàm số $y=\dfrac{-x^2+x+1}{x+1}$ có đồ thị $(C)$. Viết phương trình đường thẳng đi qua điểm $A(-1;5)$ và tiếp xúc với đồ thị $(C)$.
\loigiai
{Tập xác định $\mathscr{D}=\mathbb{R}\setminus\{-1\}$.\\
Ta có $y'=\dfrac{-x^2-2x}{(x+1)^2}$.\\
Gọi $M\left(x_0;\dfrac{-x_0^2+x_0+1}{x_0+1}\right)$  là tiếp điểm $(x_0\neq -1)$.\\
Phương trình tiếp tuyến $\Delta$  của đồ thị $(C)$ tại $M$ là $y=\dfrac{-x_0^2-2x_0}{(x_0+1)^2}\cdot (x-x_0)+\dfrac{-x_0^2+x_0+1}{x_0+1}$.\\
Vì $\Delta$ đi qua $A(-1;5)$ nên \begin{eqnarray*}
&& 5=\dfrac{-x_0^2-2x_0}{(x_0+1)^2}\cdot (-1-x_0)+\dfrac{-x_0^2+x_0+1}{x_0+1}\\
&\Leftrightarrow & 5=\dfrac{x_0^2+2x_0}{x_0+1}+\dfrac{-x_0^2+x_0+1}{x_0+1}\\
&\Leftrightarrow & 5(x_0+1)=3x_0+1\\
&\Leftrightarrow & x_0=-2.
\end{eqnarray*}
Với $x_0=2$ ta có phương trình tiếp tuyến $\Delta$ là $y=5$.\\
Vậy phương trình đường thẳng cần tìm là $y=5$.
}
\end{bt}

% nguồn: 2026-11-16_12_2D1_TL_Ham_so (nhập 16/11/2026)
% nguồn: 2026-11-16_12_2D1_TL_Ham_so (nhập 16/11/2026)
% nguồn: 2026-11-16_12_2D1_TL_Ham_so (nhập 16/11/2026)
% nguồn: 2026-11-16_12_2D1_TL_Ham_so (nhập 16/11/2026)
% nguồn: 2026-11-16_12_2D1_TL_Ham_so (nhập 16/11/2026)
% nguồn: 2026-11-16_12_2D1_TL_Ham_so (nhập 16/11/2026)
\begin{bt}%[2D1V5-6]%[ma:61463cea]
% nguon: Archive/2024/TaiLieuGiangDay/1.HamSo/1. K12 - 4in1 - Hamso/ThubaiGoc/Bai5-Dang9.tex #2822
%[Dự Án Bài Giảng 12 4 in 1, Lê Văn Toàn]
	Tìm điểm $M$ có hoành độ âm trên đồ thị $(C): y=\dfrac{1}{3}x^3-x+\dfrac{2}{3}$ sao cho tiếp tuyến tại $M$ vuông góc với đường thẳng $y=-\dfrac{1}{3}x+\dfrac{2}{3}$.
	\loigiai{
		Gọi $M\left(x_0; \dfrac{1}{3}x_0^3-x_0+\dfrac{2}{3}\right)$. \\
		Do tiếp tuyến tại $M$ vuông góc với đường thẳng $y=-\dfrac{1}{3}x+\dfrac{2}{3}$ nên ta có hệ số góc của tiếp tuyến tại $M$ là $k=3$.\\
		Ta có $y'=x^2-1$. Theo đề bài ta có phương trình $x_0^2-1=3 \Leftrightarrow x_0^2=4 \Leftrightarrow x_0=\pm 2$. \\
		Theo đề bài điểm $M$ có hoành độ âm nên $M(-2; 0)$.}
\end{bt}

% nguồn: 2026-11-16_12_2D1_TL_Ham_so (nhập 16/11/2026)
% nguồn: 2026-11-16_12_2D1_TL_Ham_so (nhập 16/11/2026)
